\documentclass[11pt]{article}
\usepackage[a4paper,margin=1.9cm]{geometry}
\usepackage[T1]{fontenc}
\usepackage{textcomp}
\usepackage[spanish,es-noshorthands]{babel}
\usepackage{amsmath,amssymb}
\usepackage{enumitem}
\usepackage{tikz}
\usetikzlibrary{calc,arrows.meta,patterns,decorations.pathmorphing}
\usepackage{xcolor}
\usepackage{multicol}
\usepackage{parskip}
\usepackage{lmodern}
\renewcommand{\familydefault}{\sfdefault}

\definecolor{unalblue}{RGB}{31,73,124}
\definecolor{unalblue2}{RGB}{70,120,180}

\newlist{opciones}{enumerate}{1}
\setlist[opciones]{label=\textbf{(\alph*)},itemsep=1pt,topsep=2pt,leftmargin=1.8em,font=\normalfont}
\setlist[enumerate,1]{label=\textbf{\arabic*.},itemsep=6pt,topsep=4pt,leftmargin=1.6em,font=\normalfont}
\setlist[enumerate,2]{label=\textbf{\roman*.},itemsep=4pt,topsep=3pt,leftmargin=1.6em,font=\normalfont}
\setlist[enumerate,3]{label=\textbf{\alph*)},itemsep=2pt,topsep=2pt,leftmargin=1.4em,font=\normalfont}

\newcommand{\vect}[2]{\begin{pmatrix}#1\\#2\end{pmatrix}}
\newcommand{\vectiii}[3]{\begin{pmatrix}#1\\#2\\#3\end{pmatrix}}

\newcommand{\examheader}{%
\noindent\begin{tikzpicture}
\node[fill=unalblue, text width=\dimexpr\linewidth-16pt\relax, inner sep=8pt, rounded corners=3pt, align=left, text=white] (hdr) {%
  \begin{minipage}[c]{0.66\linewidth}
    {\Large\bfseries \examsubject}\\[2pt]
    {\small \examtype\ \textbullet\ \textcolor{white!85!unalblue}{\examsubtitle}}
  \end{minipage}\hfill
  \begin{minipage}[c]{0.28\linewidth}
    \raggedleft{\scriptsize Escuela de Matem\'aticas\\[-1pt] Universidad Nacional\\[-1pt] de Colombia, Sede Medell\'in}
  \end{minipage}%
};
\end{tikzpicture}\\[7pt]
}
\newcommand{\examfooter}{%
  \vfill
  \rule{\linewidth}{0.4pt}\\[1pt]
  {\scriptsize\textit{Transcripci\'on digital --- MathUNAL}\hfill\textcolor{unalblue}{\textbf{\thepage}}}
}
\pagestyle{empty}

\newcommand{\examsubject}{CÁLCULO EN VARIAS VARIABLES}
\newcommand{\examtype}{Segundo Parcial}
\newcommand{\examsubtitle}{Semestre 01-2014, 12 de Mayo de 2014}

\begin{document}

\examheader

\vspace{4pt}
\noindent
\begin{minipage}[t]{0.55\linewidth}
Nombre: \dotfill\\[6pt]
Profesor: \dotfill
\end{minipage}%
\hfill
\begin{minipage}[t]{0.4\linewidth}
Doc. Identidad: \dotfill\\[6pt]
Grupo: \dotfill\\[10pt]
1. \dotfill \quad 2. \dotfill \quad 3. \dotfill \quad 4. \dotfill \quad Total: \dotfill
\end{minipage}

\vspace{6pt}
\noindent\textbf{Instrucciones.} El examen tiene una duraci\'on de 1 hora y 40 minutos. \textbf{No se admiten preguntas durante el examen.} No est\'a permitido el uso de calculadora, ni tel\'efonos m\'oviles, ni cualquier otro tipo de aparato electr\'onico. Cualquier intento de fraude ser\'a sancionado con anulaci\'on del examen. En los puntos 2., 3. y 4. justifique sus respuestas.

\begin{enumerate}
\item \textbf{(36\%)} En los numerales (i)-(iv) escoja s\'olo una de las cinco opciones que se presentan. En el numeral (v) asocie la integral dada con su correspondiente gr\'afica. No se requiere que justifique sus respuestas.
\begin{enumerate}
\item \textbf{(7\%)} Sean $f(x,y) = x^2+y^2-xy+6x$ y $E=\{(x,y): y^2+x^2 \le 8\}$. De acuerdo al m\'etodo de los multiplicadores de Lagrange, el sistema que se debe resolver para hallar los extremos de $f$ sobre la frontera de $E$ es:
\begin{opciones}
\item $\begin{cases} 2x-y+6=2\lambda y \\ 2y-x=2\lambda x \\ x^2+y^2=8 \end{cases}$
\item $\begin{cases} 2x-y=2\lambda x \\ 2y-x+6=2\lambda y \\ x^2+y^2-8=0 \end{cases}$
\item $\begin{cases} 2x-y+6=2\lambda y \\ 2y-x=2\lambda x \\ x^2+y^2=8\lambda \end{cases}$
\item $\begin{cases} 2x-y+6=2\lambda x \\ 2y-x=2\lambda y \\ x^2+y^2-8=0 \end{cases}$
\item Ninguna de las anteriores.
\end{opciones}

\item \textbf{(7\%)} Sea $W$ el s\'olido en el primer octante acotado por la superficie $z=4-x^2-y^2$ y los planos $x=0$, $y=0$ y $z=0$. La integral que calcula el volumen de $W$ es:
\begin{opciones}
\item $\int_0^{2\pi}\int_0^2\int_0^{4-r^2} r\,dz\,dr\,d\theta$.
\item $\int_0^{\pi/2}\int_0^2\int_0^{4-r^2} r\,dz\,dr\,d\theta$.
\item $\int_0^{\pi}\int_0^2\int_0^{4-r^2} r\,dz\,dr\,d\theta$.
\item $\int_0^{\pi/2}\int_0^4\int_0^{4-r} r\,dz\,dr\,d\theta$.
\item Ninguna de las anteriores.
\end{opciones}

\item \textbf{(7\%)} Sean $B=\{(x,y)\in\mathbb{R}^2: 0\le y\le x^4 \text{ y } 0\le x\le 1\}$ y $f(x,y)=e^{x^2+y^2}$. Se sabe que el valor m\'inimo absoluto de $f$ en $B$ es $1$ y que el valor m\'aximo absoluto de $f$ en $B$ es $e^2$. Entonces el valor de $\iint_B e^{x^2+y^2}\,dA$ se encuentra entre:
\begin{opciones}
\item $1$ y $e^2$.
\item $1/5$ y $e^2/5$.
\item $3/5$ y $3e^2/5$,
\item $4/5$ y $4e^2/5$
\item Ninguna de las anteriores.
\end{opciones}

\item \textbf{(7\%)} Sea $E$ el s\'olido en el primer octante limitado por el cilindro parab\'olico $z=1-x^2$ y por los planos $z=0$, $y=0$, $x=0$ y $y=4-x$. La siguiente integral \textbf{NO} calcula el volumen de $E$:
\begin{opciones}
\item $\int_0^1\int_0^{4-x}\int_0^{1-x^2} dz\,dy\,dx$.
\item $\int_0^1\int_0^{1-x^2}\int_0^{4-x} dy\,dz\,dx$.
\item $\int_0^1\int_0^{4-y}\int_0^{1-x^2} dz\,dx\,dy$.
\item $\int_0^1\int_0^{\sqrt{1-z}}\int_0^{4-x} dy\,dx\,dz$.
\item Ninguna de las anteriores. Es decir, todas las anteriores calculan el volumen de $E$.
\end{opciones}

\item \textbf{(8\%)} Asocie cada regi\'on con la integral dada

\begin{center}
\begin{tabular}{|l|c|}
\hline
\textbf{Integral} & \textbf{Gr\'afica} \\
\hline
$\int_0^1\int_0^{x/2} f(x,y)\,dy\,dx$ & \dotfill \\
\hline
$\int_0^1\int_0^{y} f(x,y)\,dx\,dy$ & \dotfill \\
\hline
$\int_0^1\int_{x/2}^{x} f(x,y)\,dy\,dx$ & \dotfill \\
\hline
$\int_0^1\int_{y/2}^{1} f(x,y)\,dx\,dy$ & \dotfill \\
\hline
\end{tabular}
\hfill
\begin{tikzpicture}[scale=1.5]
\begin{scope}[xshift=0cm]
\draw[->] (0,0) -- (1.3,0) node[right] {\scriptsize $x$};
\draw[->] (0,0) -- (0,1.3) node[above] {\scriptsize $y$};
\draw[gray!50] (0,0.5) -- (1.3,0.5); \draw[gray!50] (0,1) -- (1.3,1); \draw[gray!50] (1,0) -- (1,1.3);
\fill[gray!60] (0,0) -- (1,0.5) -- (1,1) -- (0,0) -- cycle;
\node at (0.65,-0.3) {\scriptsize (I)};
\node at (-0.15,1) {\scriptsize 1};
\node at (-0.2,0.5) {\scriptsize $\frac{1}{2}$};
\node at (1,-0.15) {\scriptsize 1};
\end{scope}
\begin{scope}[xshift=3.3cm]
\draw[->] (0,0) -- (1.3,0) node[right] {\scriptsize $x$};
\draw[->] (0,0) -- (0,1.3) node[above] {\scriptsize $y$};
\draw[gray!50] (0,1) -- (1.3,1); \draw[gray!50] (1,0) -- (1,1.3);
\fill[gray!60] (0,0) -- (1,0) -- (1,1) -- cycle;
\node at (0.65,-0.3) {\scriptsize (II)};
\node at (-0.15,1) {\scriptsize 1};
\node at (1,-0.15) {\scriptsize 1};
\end{scope}
\end{tikzpicture}

\vspace{6pt}
\begin{tikzpicture}[scale=1.5]
\begin{scope}[xshift=0cm]
\draw[->] (0,0) -- (1.3,0) node[right] {\scriptsize $x$};
\draw[->] (0,0) -- (0,1.3) node[above] {\scriptsize $y$};
\draw[gray!50] (0,1) -- (1.3,1); \draw[gray!50] (0.5,0) -- (0.5,1.3); \draw[gray!50] (1,0) -- (1,1.3);
\fill[gray!60] (0,0) -- (0.5,0) -- (1,1) -- (0,1) -- cycle;
\node at (0.65,-0.3) {\scriptsize (III)};
\node at (-0.15,1) {\scriptsize 1};
\node at (0.5,-0.15) {\scriptsize $1/2$};
\node at (1,-0.15) {\scriptsize 1};
\end{scope}
\begin{scope}[xshift=3.3cm]
\draw[->] (0,0) -- (1.3,0) node[right] {\scriptsize $x$};
\draw[->] (0,0) -- (0,1.3) node[above] {\scriptsize $y$};
\draw[gray!50] (0,0.5) -- (1.3,0.5); \draw[gray!50] (0,1) -- (1.3,1); \draw[gray!50] (1,0) -- (1,1.3);
\fill[gray!60] (0,0) -- (1,0.5) -- (0,0) -- cycle;
\node at (0.65,-0.3) {\scriptsize (IV)};
\node at (-0.15,1) {\scriptsize 1};
\node at (-0.2,0.5) {\scriptsize $\frac{1}{2}$};
\node at (1,-0.15) {\scriptsize 1};
\end{scope}
\end{tikzpicture}
\end{center}
\end{enumerate}

\vspace{6pt}
\item \textbf{(15\%)} Sea $W$ el s\'olido que est\'a dentro de la superficie $x^2+y^2+z^2=2z$ y fuera de la superficie $x^2+y^2+z^2=z$. Suponga que en todo punto de $W$ la densidad es directamente proporcional a su distancia al origen. Dibuje la regi\'on $W$ y \textbf{plantee} (sin resolver) la integral que calcula su masa en coordenadas esf\'ericas.
\end{enumerate}

\examfooter
\newpage

\examheader

\begin{enumerate}
\setcounter{enumi}{2}
\item \textbf{(24\%)} Sea $D\subseteq\mathbb{R}^2$ la regi\'on del primer cuadrante acotada por las curvas $y=0$, $y=1$, $x=\sqrt[3]{y}$ y $x=1$. Dibuje la regi\'on $D$ y calcule
\[
\iint_D \sqrt{1+x^4}\,dA.
\]
\end{enumerate}

\examfooter
\newpage

\examheader

\begin{enumerate}
\setcounter{enumi}{3}
\item \textbf{(25\%)} Sea $D$ la regi\'on del plano acotada por las rectas $y=1-x$, $y=2-x$, $y=x-1$ y $y=x+1$. Calcule
\[
\iint_D \frac{\cos(y-x)}{y+x}\,dA.
\]

Dibuje claramente la regi\'on de integraci\'on de toda integral que escriba.
\end{enumerate}

\examfooter

\end{document}
